% Generated draft: requires mathematical and source review.
\documentclass[11pt]{article}
\usepackage[margin=25mm]{geometry}
\usepackage{fontspec}
\setmainfont{Latin Modern Roman}
\usepackage{amsmath,amssymb,mathtools,unicode-math}
\setmathfont{Latin Modern Math}
\usepackage{enumitem,needspace,xurl,hyperref}
\setlength{\parindent}{0pt}
\setlength{\parskip}{5pt}
\setlength{\emergencystretch}{3em}
\newcommand{\N}{\mathbb{N}}
\newcommand{\Z}{\mathbb{Z}}
\newcommand{\Q}{\mathbb{Q}}
\newcommand{\R}{\mathbb{R}}
\newcommand{\C}{\mathbb{C}}
\newcommand{\F}{\mathbb{F}}
\newcommand{\A}{\mathbb{A}}
\newcommand{\PP}{\mathbb{P}}
\newcommand{\E}{\mathbb{E}}
\newcommand{\eps}{\varepsilon}
\newcommand{\dd}{\,\mathrm{d}}
\begin{document}
\section*{UnsolvedMath ID: 700552}
\textbf{Source problem code:} catalogue entry
\subsection{Definitions and mathematical statement}
Let \(G\) be an abstract group. Its commutator subgroup is
\[
[G,G]=\langle x^{-1}y^{-1}xy:x,y\in G\rangle.
\]
A group is residually finite if, for every \(g\ne 1\), there exist a finite group \(F\) and a homomorphism \(\varphi:G\to F\) such that \(\varphi(g)\ne 1\). An abelian group \(A\), written additively, is divisible if
\[
\forall a\in A\ \forall n\in\Z_{\geq 1}\ \exists b\in A\quad nb=a.
\]
An embedding means an injective group homomorphism.

\textbf{Unresolved terminology.} The supplied record does not define ``locally normal'' or specify the convention for ``direct product.'' These omissions are material. In particular, an unrestricted Cartesian product \(\prod_{i\in I}F_i\) and the subgroup
\[
\{(f_i)\in\prod_{i\in I}F_i:f_i=1\text{ for all but finitely many }i\}
\]
are different constructions. Ordinary residual finiteness already implies embeddability in an unrestricted Cartesian product of finite groups, which makes that interpretation problematic in conjunction with the supplied archived negative answer to part (a). The intended convention cannot be established from this record alone.

To preserve the questions without silently resolving these issues, let \(\mathcal L\) denote the class meant by the source's term ``locally normal.'' Write \(P_{\mathrm{fin}}(K)\) for the assertion that \(K\) embeds in a ``direct product'' of finite groups under the intended source convention, and write \(P_{\mathrm{cyc}}(K)\) for the assertion that \(K\) is isomorphic to a ``direct product'' of cyclic groups under that convention. The record imposes no explicit finiteness restriction on the cyclic factors. These predicates retain unresolved source terminology; the following are therefore qualified formulations requiring source clarification.

\textbf{Part (b), the remaining question in the supplied record.} Is the following assertion true?
\[
\forall G\in\mathcal L\quad
\bigl(G\text{ is residually finite}\bigr)
\Longrightarrow
\exists H\trianglelefteq G\quad
\bigl(P_{\mathrm{fin}}(H)\ \land\ G/H\text{ is divisible abelian}\bigr).
\]
No requirement that \(H\) be proper or that \(G/H\) be nontrivial is stated.

\textbf{Part (a), retained as archived context.} The original question asks whether
\[
\forall G\in\mathcal L\quad
\bigl(G\text{ is residually finite}\ \land\ P_{\mathrm{cyc}}(G/[G,G])\bigr)
\Longrightarrow P_{\mathrm{fin}}(G).
\]
The supplied archive reports a negative answer to part (a), attributed to L. A. Kurdachenko (1986). This is a supplied historical assessment, not an independently verified conclusion here. The registry labels the combined record partially solved, with part (b) retained in the main section.
\subsection{Short English statement}
Part (b) asks whether every locally normal residually finite group has a normal subgroup that embeds in a direct product of finite groups and has divisible abelian quotient. Archived part (a) asks whether such a group itself embeds in a direct product of finite groups when its abelianization is a direct product of cyclic groups; the supplied archive reports a negative answer. The meanings of ``locally normal'' and ``direct product'' require clarification before these formulations can be considered fully precise.
\subsection{Sources}
\begin{itemize}
\item Supplied local UnsolvedMath registry record: permanent ID \(700552\), exact source problem code \texttt{catalogue entry}, title \emph{Kourovka Notebook Problem 9.17}; upstream revision \texttt{d371ee917ac9647fbb06cd6fc529deec6f3dda82}.
\item \emph{The Kourovka Notebook (Unsolved Problems in Group Theory)}, 9th issue (1984), Problem 9.17, proposed by L. A. Kurdachenko. The supplied background identifies an update in the 21st issue (2026), \texttt{arXiv:1401.0300v48}. These source details are supplied attribution and have not been independently checked.
\item Supplied Archive of Solved Problems entry: ``Not always,'' attributed to L. A. Kurdachenko, \emph{Math. Notes}, volume 39 (1986), pages 273--279. The record associates this archived resolution with part (a); it supplies no article title, theorem number, or explanation of the counterexample.
\item The supplied status review is dated \texttt{2026-10-16}, records evidence as ``published,'' and assigns the canonical status \texttt{partially\_solved}. Its notes place solved parts in the archive and remaining parts in the main section. This assessment is supplied dated curation, not independently verified literature assessment.
\end{itemize}
\end{document}
